\section{Introduction}

Spinodal decomposition is a phase separation process characterized by barrierless uphill diffusion resulting in a bicontinuous microstructure morphology with periodic oscillations in composition. Since properties arise from structure, it is important to identify regions of conditions space which may be unstable to composition fluctuations and susceptible to spontaneous phase separation when designing alloys. In some cases, spinodal decomposition is undesirable and can lead to significant hardening and embrittlement of single-phase ductile alloys \cite{wang2023spinodal} or provide a a pathway for the formation of brittle, deleterious phases \cite{dasari2023crystallographic}. On the other hand, spinodal and spinodal-assisted microstructures have recently been leveraged to overcome the strength-ductility tradeoff \cite{an2021spinodal,lee2025spinodal,zhang2022refinement}, tailor magnetic properties \cite{rao2021beyond,dong2024cooperative}, and improve thermoelectric performance \cite{duan2025spinodal}. 

Regardless of whether the designer is looking to exploit or avoid spinodal decomposition, such a region must first be located within a design space that is often high-dimensional. Because exhaustive experimental characterization of these large design spaces is infeasible, computational modeling approaches which can predict microstructure stability are essential for the efficient identification of spinodal regions. Recent work to this end includes high-throughput \textit{ab initio} calculations coupled with cluster expansion to construct CALPHAD-style Gibbs energy models \cite{usanmaz2016first} and a novel procedure not requiring any functional form of the free energy in which the spinodal decomposition process is modeled as a nucleation event infinitesimally close to the spinodal boundary \cite{divilov2024priori}. While these approaches do not require any empirical data and were well-validated with experiments, they were only demonstrated on binary systems. Showing promise for spinodal design in high-dimensional composition spaces, Koneru, Kadirvel, and Wang recently developed a CALPHAD-based high-throughput microstructure stability predictive framework specifically for the multi-principal element alloy (MPEA) or high entropy alloy (HEA) design problem \cite{koneru2022high,kadirvel2022exploration}. This framework, now implemented in Pandat \cite{cao2009pandat}, facilitates gauge-invariant prediction and visualization of summaries of Gibbs energy surface geometry such as eigenvalues and eigenvectors of the Gibbs energy Hessian, spinodal temperatures, and initial composition fluctuations. 


In this work, we start with the framework of Koneru, Kadirvel, and Wang as a foundation and demonstrate a few key extensions. First, rather than using numerical derivative approximations which can become computationally expensive as the number of components gets large, we employ the recently formalized Jansson derivative method to calculate analytic Gibbs energy Hessians in a robust and efficient manner \cite{kunselman2024analytically}. Second, instead of relying on access to commercial databases which may not contain models for all alloy systems of interest, we use the recently developed PhaseForge package \cite{sariturk_2025_15730911} to automatically construct multicomponent CALPHAD models from machine learning-based interatomic potentials (MLIP). Third, we include an estimate for the elastic energy contributions due to coherency strains which penalize phase separation and effectively shrink the miscibility gap and spinodal region. Fourth, instead of using a collection of constant core component (a type of pseudoternary) diagrams to visualize results, we employ affine projections \cite{vela2025visualizing} to summarize high-dimensional data on a single two-dimensional plot. 

Lastly, the entire workflow is packaged in an open-source tool available through the MaterialsFramework interface \cite{sariturk_2025_15731044}. Through this framework, modelers have whitebox access to the CALPHAD model, and as demonstrated in \cite{kunselman2025construction}, they can further refine parameters to any available computational or experimental data through ESPEI \cite{bocklund2019espei}. The remainder of this work will dive into the theory and methodology of determining microstructure stability from the Gibbs energy description of a phase and then demonstrate our framework to investigate the stability of the BCC phase in the Hf-Nb-Ti-V system.
